% Part: first-order-logic
% Chapter: completeness
% Section: introduction

\documentclass[../../../include/open-logic-section]{subfiles}

\begin{document}

\iftag{FOL}
      {\olfileid[fr]{fol}{com}{int}}
      {\olfileid[fr]{pl}{com}{int}}

\olsection{Introduction}

Le théorème de complétude est l'un des résultats fondamentaux de la
logique. Il admet deux formulations dont nous démontrerons
l'équivalence. Dans la première, il exprime un lien essentiel entre
la conséquence sémantique et notre système de !!{derivation}s : si
\iftag{FOL}{!!a{sentence}}{!!a{formula}}~$!A$ est une conséquence \iftag{FOL}{d'!!{sentence}s}{de !!{formula}s}~$\Gamma$,
il existe aussi !!a{derivation} établissant
$\Gamma\Proves !A$. Le système de !!{derivation}s est donc aussi
puissant qu'il peut l'être sans démontrer de \iftag{FOL}{énoncés}{formules} qui ne découlent
pas des hypothèses.

Selon la seconde formulation, le théorème affirme l'existence de
modèles : tout ensemble cohérent \iftag{FOL}{d'!!{sentence}s}{de !!{formula}s} est satisfaisable.
La cohérence est une notion de théorie de la démonstration.
Elle signifie que l'absurde ne se dérive pas à partir de~$\Gamma$.
Pourquoi cette propriété garantirait-elle qu'il existe
\iftag{FOL}{!!a{structure}~$\Struct{M}$ telle que
$\Sat{M}{\Gamma}$}{!!a{valuation}~$\pAssign{v}$ telle que
$\pSat{v}{\Gamma}$} ? Avant la première démonstration du théorème
de complétude, et même avant l'apparition des systèmes de
!!{derivation}s que nous utilisons aujourd'hui, le mathématicien
allemand David Hilbert défendait l'idée que la cohérence des théories
mathématiques garantit l'existence des objets dont elles parlent.
Dans une lettre à Gottlob Frege, il l'exprima ainsi :
\begin{quote}
  Si les axiomes posés arbitrairement ne se contredisent pas, avec
  toutes leurs conséquences, alors ils sont vrais et les choses
  définies par ces axiomes existent. C'est pour moi le critère de la
  vérité et de l'existence.
\end{quote}
Frege s'y opposait fermement. La seconde formulation du théorème de
complétude montre que Hilbert avait raison au moins en ce sens : si
les axiomes sont cohérents, il existe \emph{au moins une}
\iftag{FOL}{!!{structure}}{!!{valuation}} qui les rend tous vrais.

Ce ne sont pas les seules raisons de l'importance du théorème de
complétude ou, plus exactement, de sa démonstration. Celle-ci a
plusieurs conséquences importantes, dont certaines seront étudiées
séparément. Par exemple, toute !!{derivation} qui établit
$\Gamma\Proves !A$ n'utilise qu'un nombre fini \iftag{FOL}{d'!!{sentence}s}{de !!{formula}s}
de~$\Gamma$. Par complétude, puis par correction,
si $!A$ est une conséquence de~$\Gamma$, elle est déjà une
conséquence d'un sous-ensemble fini de~$\Gamma$. C'est la
\emph{compacité}. Grâce à l'équivalence entre cohérence et
satisfaisabilité, on peut aussi dire que $\Gamma$ est cohérent
dès que tous ses sous-ensembles finis le sont.

Bien que le théorème de compacité découle ainsi de la complétude par
l'intermédiaire des !!{derivation}s, on peut aussi l'établir
directement au moyen de la \emph{démonstration} du théorème de
complétude. Cette démonstration part d'un ensemble cohérent
\iftag{FOL}{d'!!{sentence}s}{de !!{formula}s} et construit
\iftag{FOL}{une !!{structure}}{une !!{valuation}} qui le satisfait.
Presque la même construction établit directement la compacité en
partant d'un ensemble \iftag{FOL}{d'!!{sentence}s}{de !!{formula}s} \emph{finiment satisfaisable}.
\iftag{FOL}{La construction a d'autres conséquences. Elle montre par
exemple que tout ensemble satisfaisable d'!!{sentence}s possède un
modèle fini ou infini dénombrable dans le langage au plus dénombrable
fixé ici. Ce résultat est appelé le théorème de
L\"owenheim--Skolem. Plus généralement, construire des
!!{structure}s à partir d'ensembles d'!!{sentence}s est une méthode
fréquente en logique et parfois même en philosophie.}{}

\end{document}
